\chapter*{Abstract}
\addcontentsline{toc}{chapter}{Abstract}

Large language models deployed in serverless inference environments are frequently evicted from GPU memory when idle and must reload their parameters on demand. This checkpoint loading step dominates cold-start latency for large models. Although modern NVMe devices provide substantial peak throughput, effective loading performance frequently falls short of hardware limits. This discrepancy suggests that loading performance is constrained less by storage bandwidth and more by how read requests are submitted to the operating system.

This project presents TensorStore, a checkpoint loading framework that supports multiple storage backends and checkpoint formats, and evaluates them under controlled and reproducible conditions. The framework provides SafeTensors compatibility and a partitioned ServerlessLLM-style format, with backends spanning synchronous POSIX I/O, memory-mapped access, and Linux io\_uring. Python bindings enable integration testing within representative serving workflows.

A systematic evaluation quantifies the effect of backend choice, checkpoint layout, and integration layer on loading performance. Results show that synchronous I/O underutilises NVMe hardware due to blocking submission and shallow queue depth, while io\_uring achieves higher throughput and read IOPS by sustaining deeper queues. The advantage is largest for partitioned checkpoint layouts. In the vLLM integration, backend-level improvements propagate partially to end-to-end cold-start latency, as loading is one component of initialisation rather than the whole of it. Findings indicate that checkpoint loading in serverless LLM systems is often submission-bound rather than bandwidth-bound, and that async I/O mechanisms deserve wider attention in inference storage design.
